\section{OBJECTIVE AND SOURCE REVIEW}
\textbf{Erd\H{o}s \#222; independent attempt, 2026-09-23.}
The problem asks for good bounds on gaps between successive sums of two integer squares. I read the full \texttt{222.tex}. Its quarter-power construction is correct and classical; calling it sharp does not establish an optimal exponent. I examine the complementary constructive lower-bound mechanism explicitly.
\section{WORK: A CERTIFIABLE GAP OF ANY FIXED LENGTH}
Choose distinct primes $q_1,\ldots,q_H\equiv3\pmod4$. The Chinese remainder theorem supplies a residue $X$ modulo $Q=\prod_{j=1}^Hq_j^2$ satisfying
\[
X+j\equiv q_j\pmod {q_j^2}\quad(1\le j\le H).
\]
Every positive representative of this residue gives a block $X+1,\ldots,X+H$ with no sum of two squares. Indeed, if $q\equiv3\pmod4$ divides $a^2+b^2$, then $q\mid a,b$: otherwise $-1$ would be a square modulo $q$, contradicting Fermat's theorem. Thus such a prime cannot divide a sum of squares to exact exponent one, whereas our congruence forces precisely that exponent.
Taking $0\le X<Q$, a gap between consecutive representable integers surrounding this block has length at least $H+1$ and begins below $Q$. Thus, with the maximal-gap convention of the source,
\[
g(Q)\ge H+1.
\]
There are arbitrarily many primes $3\pmod4$: given a finite list, a prime factor $3\pmod4$ of $4\prod q_j-1$ is outside it. So this is a completely constructive unbounded-gap proof, including the location bound in terms of its selected primes.
\section{ADVERSARIAL VERIFICATION AND REMAINING GAP}
One must impose residues modulo $q_j^2$, not merely modulo $q_j$, to exclude an even valuation. The block endpoint need not itself be a sum of squares; surrounding it by the nearest two representable integers only increases the gap. The large product $Q$ makes this weaker than the known logarithmic lower bounds recorded in the official entry. It is a verified certificate framework, not a new record. Neither it nor the supplied quarter-power upper bound determines the true order of the largest gaps.
\section{SOURCES CHECKED}
Complete local source: \texttt{gpt\_pro\_5.4/222.tex}. Official indexed remarks checked at \url{https://www.erdosproblems.com/search_bib/Er57?sources_only=1}, entry 222, 2026-09-23. No advanced external theorem is used in the CRT argument.
\section{FINAL}
\textbf{UNRESOLVED.} The constructive lower obstruction is quantified, with no verified improvement to the best general bounds.
\section{COMPLETION ESTIMATE}
Subjective progress toward determining gap growth.
\textbf{COMPLETION: 10\%}
